% Supplement: https://eel.is/c++draft/basic.fundamental#2
% Collision reference: https://cp-algorithms.com/string/string-hashing.html
\subsection{字符串自然溢出哈希}
\label{algo:string-hash-overflow}
\AlgorithmInfo
    {省去显式取模，用 64 位无符号整数维护子串哈希。}
    {预处理时间、空间均为 $O(n)$；每次区间查询或比较 $O(1)$。}
    {使用 \InlineCode{uint64_t}；查询为从 1 开始的闭区间，空区间返回 0。}
\subsubsection*{原理}
与单哈希使用相同的前缀递推和区间消去公式，只把模数换成 $2^{64}$：
\[
    h_i=h_{i-1}B+c_i,\qquad
    H(l,r)=h_r-h_{l-1}p_{r-l+1}\pmod{2^{64}}.
\]
\InlineCode{uint64_t} 的加、减、乘运算会自动保留结果的低 64 位，等价于模 $2^{64}$；减法回绕也能得到正确余数。因此不需要写 \texttt{\%mod} 或在减法后补模数，代码更简短。
\subsubsection*{使用教程}
\InlineCode{Hash64} 的接口与单哈希一致。先初始化字符串，再查询哈希或比较两个子串；保存结果时用 \InlineCode{auto} 或 \InlineCode{uint64_t}。
\begin{lstlisting}[style=mtycpp,numbers=none]
Hash64 ha;
ha.init("ababa");
auto h=ha.query(1,3);
bool ok=ha.same(1,3,3,5);
\end{lstlisting}
不同字符串分别初始化对象，在区间长度相同的前提下比较 \InlineCode{query} 返回值。\InlineCode{init} 可重复调用；非空区间须满足 $1\le l\le r\le n$。
\needspace{7\baselineskip}
\subsubsection*{模板代码}
\CodeFile{code/strings/hash_overflow.cpp}
\subsubsection*{需要注意}
前缀、幂与返回值都应保持无符号 64 位，不能用 \InlineCode{long long} 的有符号溢出来代替，也不能用 \InlineCode{1ULL<<64} 表示模数。自然溢出同样存在碰撞，而且可构造针对模 $2^{64}$ 的碰撞；位数较多不代表一定比单模数更可靠，对抗数据优先使用双模数哈希。
